\scanpage{73}

% top center and right
\mspage{- 139 -}\hfill IX

% upper, continuation
définir l'isomorphisme canonique (11.6.2).

% upper, inductive limit
Ecrivant $M_\ell\otimes D_\ell$ (resp. $\check{M}_\ell\otimes D_\ell$) comme limite inductive des $M_n$ (resp. des $\check{M}'_n$), où $n$ parcourt les puissances de $\ell$, il revient au même de définir (11.6.2), où des isomorphismes canoniques

% upper, (11.6.4)
\[
{}_n\Phi_{\unsure{\mathrm{o}|0}}\simeq\operatorname{Ker}(\widetilde{u}_n:M_n\longrightarrow\check{M}'_n)\quad,\quad n=\ell^{r+1}\quad,\tag*{(11.6.4)}
\]

% upper middle, after (11.6.4)
comptatibles avec les morphismes de transition évidents, pour $n$ variable\handwritten{\unsure{!}} (où $\widetilde{u}_n$ désigne maintenant, par abus de notations, l'homomorphisme déduit de $\widetilde{u}_\ell:M_\ell\longrightarrow\check{M}'_\ell$ par réduction mod. $n$).

% middle, canonical filtration
Pour ceci, considérons le schéma en groupes $ {}_nA$ sur $S$, qui est plat et quasi-fini sur $S$ (2.2.1), grâce à l'hypothèse de réduction semi-stable. On a une filtration canonique

% right margin at filtration paragraph
\margin{\handwritten{$\unsure{\lrcorner}$ \unsure{15}}}

% middle, (11.6.5), three right-hand hooks
\[
{}_nA\mathrel{\relbar\joinrel\rhook}({}_nA^{\mathrm{o}})^f\mathrel{\relbar\joinrel\rhook}({}_nA^{\mathrm{o}})^t\mathrel{\relbar\joinrel\rhook}0\quad,\tag*{(11.6.5)}
\]

% middle lower, exponents and completeness
où l'exposant $f$ (resp. l'exposant $t$) désigne la ``partie fixe'' (2.2.3) (resp. la ``partie torique'' (5.1)) du schéma en groupe qu'il affecte ; pour la définition de la partie torique dans 5.1, on notera qu'on avait supposé $S$ complet, ce qui est loisible ici (11.0). (On a d'ailleurs vu dans 6.1 que la construction se généralise au cas $S$ hensélien, si on y tient \ldots) Soit

% lower, (11.6.6)
\[
{}_n\Psi={}_nA/({}_nA^{\mathrm{o}})^f\quad,\tag*{(11.6.6)}
\]

% lower, generic fibre and references
c'est un schéma en groupes plat, séparé et quasi-fini sur $S$ (SGA 3 $\mathrm{VI}_A$ 3.2 et $\mathrm{VI}_B$ 9.2) de fibre générique

% lower, generic-fibre equation
\[
{}_n\Psi_K=({}_nA_K)/({}_nA_K)^f\quad,
\]

% bottom strip, faint numerals
\handwritten{\unsure{457}\qquad\unsure{13}}

% bottom center, printed label
457

% lower right
% archive sheet number: 72
